\documentclass[11pt,a4paper]{article}
\usepackage[T1]{fontenc}
\usepackage[spanish,es-nodecimaldot]{babel}
\usepackage[margin=2.2cm]{geometry}
\usepackage{tabularx,booktabs}
\usepackage{csquotes}
\usepackage[backend=biber,style=ieee]{biblatex}
\usepackage[hidelinks]{hyperref}
\addbibresource{referencias.bib}
\setlength{\parindent}{0pt}
\setlength{\parskip}{4pt}
\renewcommand{\arraystretch}{1.14}

\begin{document}
{\large\bfseries FORM-AC-01\par}
{\LARGE\bfseries Caída de X del 10 de marzo de 2025\par}
{Melanie Michelle Muñoz Vera\par}

\section*{Descripción del incidente}
El 10 de marzo de 2025, X (antes Twitter) sufrió varias interrupciones que dificultaron el acceso a su aplicación móvil y a su sitio web. Cisco ThousandEyes detectó errores al establecer conexiones TCP en varias regiones y observó una pérdida de paquetes compatible con un ataque de denegación de servicio distribuido (DDoS) \cite{hicks-x-2025}. Ese mismo día, Elon Musk atribuyó la interrupción a un ciberataque \cite{chapman-ortutay-x-2025}. El servicio presentó fallos intermitentes durante varias horas. Las fuentes consultadas no permiten fijar una hora de recuperación global \cite[p.~15]{csc-threat-update-2025}.

\section*{Línea de tiempo}
Horas aproximadas del 10 de marzo en UTC. Las horas de Associated Press se convirtieron desde EDT (UTC$-4$).

\begin{tabularx}{\linewidth}{@{}l>{\raggedright\arraybackslash}X@{}}
\toprule
\textbf{Hora (UTC)} & \textbf{Qué ocurrió} \\
\midrule
09:45 & ThousandEyes empezó a detectar una caída de disponibilidad y errores de conexión TCP en varias regiones \cite{hicks-x-2025}. \\
10:00 & Downdetector registró un primer pico de avisos de problemas de acceso, según Associated Press \cite{chapman-ortutay-x-2025}. \\
14:00 & Se produjo un segundo pico, con más de 40\,000 reportes \cite{chapman-ortutay-x-2025}. \\
16:00 & Comenzó otro episodio de interrupción sostenida que duró al menos una hora, con mayor afectación en las costas de EE.\,UU., según Associated Press \cite{chapman-ortutay-x-2025}. \\
\bottomrule
\end{tabularx}

\section*{Impacto cuantificado}
Del conjunto de problemas registrados por Downdetector, el 56\,\% correspondió a la aplicación móvil y el 33\,\% al sitio web. El pico superó los 40\,000 reportes. Esa cifra mide avisos enviados al monitor y deja sin determinar el total mundial de personas afectadas. La duración mínima de una hora corresponde al episodio de las 16:00 UTC \cite{chapman-ortutay-x-2025}.

\printbibliography[title={Referencias}]
\end{document}